\chapter{重正化、匹配与胶子 TMD 的系统构造}

\section{从共线分布到横动量依赖分布}

共线胶子分布描述纵向动量份额，而 TMD-PDF 同时保留横向分离。采用伴随表示
Wilson 线的非极化胶子定义可以写成
\begin{equation}
  x\,g(x,\mu^2)
  =\frac{1}{2\pi P^+}\int \dd\xi^- e^{-ixP^+\xi^-}
  \langle P|F^{+i,a}(\xi^-)W^{ab}_{\mathrm{adj}}(\xi^-,0)
  F^{+i,b}(0)|P\rangle_{\mu},
  \label{eq:b10-lightcone-pdf}
\end{equation}
其中 $P^+$ 是光锥正方向动量，$i=1,2$ 是横向指标，$W_{\mathrm{adj}}$ 是
保证规范不变性的路径有序 Wilson 线。式~\eqref{eq:b10-lightcone-pdf} 的
物理意义是沿光锥分离的胶子场强关联；它定义在 Minkowski 光锥上，因此不能
直接替换为欧氏格点上的等时关联。若在光锥规范中令 $A^+=0$，Wilson 线化为
单位矩阵，但规范无关的物理结论不能依赖这一选择。

TMD 定义在同一纵向结构上增加横向位移 $\bperp$ 与 staple 规范链：
\begin{equation}
  x\,g(x,\bperp;\mu,\zeta)
  =\frac{1}{2\pi P^+}\int \dd\xi^- e^{-ixP^+\xi^-}
  \langle P|F^{+i}(\xi^-+\bperp)W_{\mathrm{staple}}
  F^{+i}(0)|P\rangle_{\mu,\zeta}.
  \label{eq:b10-tmd-def}
\end{equation}
这里 $\zeta$ 是 Collins--Soper 快度标度。与式~\eqref{eq:b10-lightcone-pdf}
相比，式~\eqref{eq:b10-tmd-def} 多出三个不可省略的结构：横向分离、同几何
软减除以及 Wilson 线方向造成的过程依赖。只有当 $q_\perp\ll Q$ 且横动量
因子化成立时，TMD 才能作为截面中的独立非微扰输入；把任意大 $q_\perp$ 的
数据写成 TMD 会越过因子化窗口。

在格点侧，常用准 TMD 矩阵元 $\tilde g(x,\bperp,\Pz)$ 逼近式
\eqref{eq:b10-tmd-def} 的光锥对象。其大动量展开应理解为
\begin{equation}
  \tilde g(x,\Pz)
  =\int_{-1}^{1}\frac{\dd y}{|y|}
  C\!\left(\frac{x}{y},\frac{\mu}{|y|\Pz}\right)g(y,\mu)
  +\Ocal\!\left(\frac{\LQCD^2}{\Pz^2},
  \frac{\LQCD^2}{x^2\Pz^2}\right),
  \label{eq:b10-lamet}
\end{equation}
其中 $C$ 是微扰匹配核，后两项分别是普通高扭度/靶质量修正和小 $x$ 增强。
该式是有效场论展开，不是对所有 $\Pz$ 的恒等式。实际分析至少要求
$\Pz\gtrsim1.5~\mathrm{GeV}$ 且 $x\Pz\gg\LQCD$；小 $x$ 端若
$x\Pz$ 落入非微扰区，匹配核与 Landau 极点问题不能被统计误差掩盖。

\section{欧氏算符、投影与 Wilson 线几何}

令 Wilson 线沿 $z$ 方向，横向方向为 $i,j\in\{x,y\}$。胶子双局域算符可写为
\begin{equation}
  M^{\mu\lambda;\nu\rho}(z,\bperp)
  =F^{\mu\lambda}(z\hat z+\bperp)U_{\mathrm{adj}}
  F^{\nu\rho}(0),
  \label{eq:b10-bilocal}
\end{equation}
其中 $U_{\mathrm{adj}}$ 是连接两端场强的伴随表示路径。非极化通道采用的
可乘组合是
\begin{equation}
  \Ocal_U(z,\bperp)
  =M^{ti;ti}+M^{tj;tj}-2M^{ij;ij},
  \qquad
  \Ocal_{U,\mathrm{even}}=\Ocal_U(+\!z)+\Ocal_U(-\!z).
  \label{eq:b10-unpolarized}
\end{equation}
若不除以 2，$z\to0$ 时奇数方向组合恒为零，可作为实现与数据布局的直接检查。
螺旋度通道的理论目标是
\begin{equation}
  \Ocal_{H,\mathrm{odd}}
  =\bigl(M^{ti;ti}_H+M^{tj;tj}_H-2M^{ij;ij}_H\bigr)_{+z}
  -\bigl(\cdots\bigr)_{-z}.
  \label{eq:b10-helicity}
\end{equation}
注意，存储格式中常见的 $M^{ti;ti}+M^{tj;tj}+2M^{ij;ij}$ 是诊断组合，
不能替代式~\eqref{eq:b10-helicity} 的物理螺旋度投影。

TMD staple 的格点几何由纵臂、横向翼和反向纵臂组成。若纵臂长度为 $L$，
并要求 $L$ 与横向位移保持在格点体积内，则 $\bperp\to0$ 时 staple 应连续
退化为直线 Wilson 线。这个极限提供三方面交叉检查：共线与 TMD 的公共归一、
$z=0$ 的局域性以及 even/odd 投影的一致性。若路径绕周期边界、纵臂与横翼
共线，或两端使用的方向指标不一致，则所有后续“规范不变”检查都不再对应同一
算符。

\section{三种尺度：UV、快度与动量}

胶子 TMD 同时含 $\mu$、$\zeta$ 与 $\Pz$。三者不能相互替代：$\mu$ 是 UV
重整化标度，$\zeta$ 是快度演化标度，$\Pz$ 是构造准光锥关联的动量。软函数
与 Collins--Soper 核的配对结构为
\begin{align}
  \frac{\dd}{\dd\ln\sqrt{\zeta}}\ln F(x,\bperp;\mu,\zeta)
  &=K(\bperp,\mu),
  \label{eq:b10-cs-evolution}\\
  \frac{\dd}{\dd\ln\mu}K(\bperp,\mu)
  &=-\gamma_\zeta(\alpha_s),
  \label{eq:b10-cs-mu}
\end{align}
其中 $F$ 是选定软减除方案后的 TMD，$K$ 是 Collins--Soper 核，
$\gamma_\zeta$ 是其 UV 反常量纲。式~\eqref{eq:b10-cs-evolution} 只在软函数、
TMD 和硬核使用同一快度方案时成立。若把一个方案的软因子除以另一个方案的
裸 TMD，或把 $\mu$ 演化与 $\zeta$ 演化混成一个拟合常数，可积性条件会被
人为破坏。

\section{梯度流重整化与小流时间展开}

有限流时间 $\tau$ 的 flowed 算符在连续极限中具有受控的 UV 性质，但它定义
的是带尺度 $\sqrt{\tau}$ 的新方案，并不自动等于 $\MSbar$。小流时间展开
写成
\begin{equation}
  \Ocal_{\mathrm{GF}}(\tau,z)
  =\sum_k C_k(\tau,\mu,z)\Ocal_k^{\MSbar}(z,\mu)
  +\Ocal(\tau\LQCD^2,\tau/z^2).
  \label{eq:b10-sfte}
\end{equation}
对 Wilson 线算符，线性自能可以组织为端点系数与指数线因子。按本批次来源的
一圈约定，
\begin{align}
  c_{\parallel\perp}(\tau,\mu)
  &=1+\Ocal(\alpha_s^2),
  \label{eq:b10-cpara}\\
  c_{\perp\perp}(\tau,\mu)
  &=1+\frac{\alpha_s(\mu)C_A}{4\pi}L_\tau(\mu)
  +\Ocal(\alpha_s^2),
  \label{eq:b10-cperp}\\
  \delta m_A(\tau)
  &=-\frac{\alpha_s(\mu)C_A}{4\pi}\frac{\sqrt{2\pi}}{\sqrt{\tau}}
  +\Ocal(\alpha_s^2),
  \label{eq:b10-dm}\\
  L_\tau(\mu)&=\ln\!\left(2\mu^2\tau e^{\gamma_E}\right).
  \label{eq:b10-Ltau}
\end{align}
于是双局域算符的一圈转换可概括为
\begin{equation}
  \Ocal_{ij}^{\MSbar}(z,\mu)
  =\frac{e^{-\delta m_A|z|}}{c_i c_j}
  \Ocal_{ij}^{\mathrm{GF},R}(\tau,z)+\Ocal(\tau),
  \qquad
  \begin{cases}
    c_i=c_j=c_{\perp\perp}, & U,\\
    \{c_i,c_j\}=\{c_{\parallel\perp},c_{\perp\perp}\}, & H.
  \end{cases}
  \label{eq:b10-endpoint}
\end{equation}
这里 $|z|$ 必须以 $\mathrm{GeV}^{-1}$ 计；若从格距单位换算，需显式乘以
$a[\mathrm{fm}]\times5.0677$。式~\eqref{eq:b10-endpoint} 只给出一圈端点
转换，不包含软函数、快度演化或连续极限。尤其不能把
\texttt{light\_cone\_matching} 的字符串标签当作匹配已经完成。

SFTE 的适用窗口不是“流时间越大越好”，而应同时满足
\begin{equation}
  \frac{a^2}{\tau}\ll1,\qquad
  \sqrt{8\tau}\ll|z|,\qquad
  \tau\Pz^2\ll1,\qquad
  \tau\LQCD^2\ll1.
  \label{eq:b10-sfte-window}
\end{equation}
第一项控制格点 UV 与流尺度的分离，第二项保证 Wilson 线仍远长于流化半径，
后两项要求流尺度低于强子与动量标度。若为了缩小误差而把大 $\tau$ 点纳入
线性外推，可能把 $\Ocal(\tau)$ 截断变成了中心值漂移；因此必须同时报告小流
窗口、大流诊断和窗口敏感性。

\section{非微扰重整化：自重整化、hybrid 与 RI/MOM}

自重整化利用零动量矩阵元的短距微扰形状和长距非微扰参数化。其思想可写成
\begin{equation}
  h_B(z,\Pz)\sim e^{-\delta m|z|}\,
  Z_{\mathrm{MS}}(z,\mu)\,B(z;\boldsymbol p),
  \label{eq:b10-zr-sketch}
\end{equation}
其中 $Z_{\mathrm{MS}}$ 在某些小 $z$ 节点由 NLO 匹配控制，$B$ 由大 $z$
行为拟合。具体实现采用含 $z/a$、$\ln(a\LQCD)$、$(m_0+m_2a^2)z$ 以及
$fa+f_2a^2$ 项的全局参数化，并通过逐 bootstrap/jackknife 样本重拟合传播
不确定度。该方案的优势是不依赖 Landau 规范固定；失效点是短距窗口、大 $z$
尾部模型和流时间不能同时受控时，参数化会以“吸收系统误差”的方式制造稳定
假象。

hybrid 方案在过渡尺度 $z_s$ 两侧拼接
\begin{equation}
  h_R(z,\Pz)=
  \begin{cases}
    \dfrac{h_B(z,\Pz)}{h_B(z,\Pz=0)}, & |z|\le z_s,\\[1.1em]
    \dfrac{h_B(z,\Pz)}{Z_R(z)}\,\eta_s, & |z|>z_s,
  \end{cases}
  \qquad
  \eta_s=\frac{Z_R(z_s)}{h_B(z_s,\Pz=0)}.
  \label{eq:b10-hybrid}
\end{equation}
短距比值消去与 $z$ 无关的 UV 因子和多数量子修正，长距 $Z_R$ 保留非微扰
信息。通常取 $z_s\sim0.2$--$0.4~\mathrm{fm}$，但 $z_s$ 是方案参数，不能
按点挑选“误差最小”的过渡位置。式~\eqref{eq:b10-hybrid} 在 $z_s$ 处连续
只说明拼接代数一致，不代表短距微扰和长距模型同时可信。

RI/MOM 通过 Landau 规范下的离壳夸克 Green 函数定义 $Z_O$，在窗口
$\LQCD\ll\mu_R\ll1/a$ 内把非微扰结果转换到 $\MSbar$。对非局域 Wilson
线，$z$ 与 $\mu_R$ 可比拟时幂次修正会污染匹配，因此通常只适合
$z\lesssim0.3~\mathrm{fm}$。B10 收录的 2022 年学位论文摘要报告：在约
$0.032$--$0.121~\mathrm{fm}$ 的多个格距和三种费米子作用量上，RI/MOM
重整化后仍可能残留 Wilson 线线性发散，且随格距减小而恶化。这个结论仅来自
摘要级阅读，本批次没有核对论文正文的投影、窗口和统计误差，不能单独据此
宣判现代 ratio、hybrid 或梯度流方案。

两种主线必须保持互斥：
\begin{align}
  \text{GF 链:}\quad
  &h^{\mathrm{GF,bare}}\longrightarrow h^{\MSbar}\longrightarrow
  \text{Fourier/matching},
  \label{eq:b10-chain-gf}\\
  \text{hybrid 链:}\quad
  &h_B\longrightarrow h_R^{(Z_R,\mathrm{hybrid})}\longrightarrow
  \text{Fourier/matching}.
  \label{eq:b10-chain-hybrid}
\end{align}
两者的输入可以是同一裸矩阵元，但输出方案不同。把 $\delta m_A$、比值因子、
$Z_R$ 或 Wilson 圈减除任意相乘，通常不及于任何定义良好的方案；分析 metadata
必须明确记录实际使用了哪一条链。

\section{Fourier 重构、NLO 卷积与软减除}

共线胶子准 PDF 的 sin 变换为
\begin{equation}
  \tilde g(x,\Pz)
  =\frac{2\Pz}{x}\int_0^{z_{\max}}\dd z\,
  h_R(z,\Pz)\sin(x\Pz z).
  \label{eq:b10-sin}
\end{equation}
准 TMD 则使用复相位与 Hermitian 延拓：
\begin{equation}
  x\,\tilde g(x,\bperp,\Pz)
  =\frac{1}{\mathcal N}\int\frac{\dd z}{2\pi\Pz}
  e^{-ix\Pz z}h_R(z,\bperp,\Pz),
  \qquad
  \mathcal N=\frac{\Pz^2}{P_t^2},
  \label{eq:b10-quasi-tmd}
\end{equation}
其中 $P_t=\sqrt{M_N^2+\Pz^2}$。在正半轴数据上采用
$h_R(-z)=h_R(z)^*$，等价于余弦乘实部、正弦乘虚部的组合。两种 Fourier
表示在有限 $z_{\max}$ 下都引入低通卷积与 Gibbs 振荡；$z_{\max}\to\infty$
并不是数值上可实现的极限，必须用大 $z$ 尾部模型或正则化估计缺失贡献。

NLO 胶子匹配的离散实现可组织为
\begin{equation}
  Z_{ij}=\delta_{ij}+\frac{\alpha_s C_A}{2\pi}
  \frac{x_i}{|x_i|}\Delta x
  \left[\frac{g_{ij}}{y_j}
  -\delta_{ij}y_i\sum_k\frac{g_{ik}}{y_k^2}\right],
  \qquad
  h^{\mathrm{PDF}}=Z^{-1}h_R^0.
  \label{eq:b10-zij}
\end{equation}
核 $g_{ij}$ 分为 $\xi=x_i/y_j<0$、$0<\xi<1$ 与 $\xi>1$ 三区，并在
$\xi=1$ 用 plus-distribution 或主值处理。矩阵求逆不是数学装饰：它同时
恢复端点 $\delta(1-\xi)$ 与软胶子求和规则。若只保留正半轴而丢掉 $\xi<0$
支路，会出现不随网格减小而消失的归一偏移；若 x 网格含精确零点，则必须停止
计算，而不是用任意有限值替代主值约定。

TMD 混合匹配把硬核、软减除和 CS 演化放在同一式中：
\begin{equation}
  x\tilde g(x,\bperp,\Pz)
  =\int\frac{\dd y}{|y|}
  C^{\mathrm{hybrid}}\!\left(\frac{x}{y},\lambda_s,
  \frac{\mu}{y\Pz}\right)
  \frac{y\,g(y,\bperp)}{\sqrt{S_I(\bperp,\mu)}}\,
  \exp\!\left[\frac12K(\bperp,\mu)
  \ln\frac{\zeta_z}{\zeta}\right].
  \label{eq:b10-tmd-match}
\end{equation}
式~\eqref{eq:b10-tmd-match} 的因子顺序对应物理层次，而不是任意约定：
$C^{\mathrm{hybrid}}$ 是微扰硬核，$S_I$ 是内禀软函数，$K$ 把参考快度演化
到目标 $\zeta$。只有三者来自同一颜色表示、同一 Wilson 线几何和同一快度
方案时，乘积才是方案无关的物理 TMD。工程上的推荐顺序是“归一、软减除、
CS 演化、NLO 卷积、联合外推”，并保存每一步的 metadata。

在短距区域，CS 核的一圈结构可写为
\begin{equation}
  K(\bperp,\mu)
  =-\frac{\alpha_s(\mu)C_A}{\pi}
  \ln\frac{\bperp^2\mu^2}{4}+\Ocal(\alpha_s^2),
  \label{eq:b10-cs-one-loop}
\end{equation}
常数项依赖方案。格点比值提取则使用两个动量：
\begin{equation}
  K(\bperp)\simeq
  \frac{\ln\!\left[h_R(z_0,\bperp,P_{z1})/
  h_R(z_0,\bperp,P_{z2})\right]}{\ln(P_{z1}/P_{z2})}.
  \label{eq:b10-cs-ratio}
\end{equation}
式~\eqref{eq:b10-cs-ratio} 只有在 $z_0$ 位于共同短距窗口、两个动量都满足
因子化且硬核差异受控时才有意义；否则它只是有效斜率，不是普适 CS 核。

\section{尚未闭合的归一化条件}

连续场强与代码 Clover 的场归一化之间可能存在树图因子 $\kappa_F$。若代码
输出没有显式包含 $g^2$，双局域算符应相应的 $\kappa_F^2$ 或等价转换吸收。
在本批次来源中，$\kappa_F$ 的微扰确定被明确列为未闭合链，而当前实现按
$\kappa_F=1$ 记录。因此，式~\eqref{eq:b10-endpoint} 的系数只有在场强归一、
Clover 投影和 Wilson 线表示全部对齐后才可用于物理比较。这个缺口不能由
ISO 误差或统计 bootstrap 自动修复。

\chapter{从组态到物理胶子 TMD 的工作流、验证与历史}

\section{端到端流程与停止条件}

一个可审计的格点胶子 TMD 流程应把数学对象、数据对象和物理结论分开。表
\ref{tab:b10-workflow} 给出最小流程；任一阶段的停止条件未满足，就只能保留
中间产物，不能把后续拟合结果称为物理 TMD。

\begin{longtable}{@{}p{0.06\textwidth}p{0.30\textwidth}p{0.50\textwidth}@{}}
\caption{B10 建议纳入第 18 章的端到端工作流。}\label{tab:b10-workflow}\\
\toprule
步 & 对象 & 进入下一步的条件 \\
\midrule
\endfirsthead
\toprule
步 & 对象 & 进入下一步的条件 \\
\midrule
\endhead
1 & ILDG thin links、几何与元数据 & 链接近 SU(3)、行列式、体积与边界语义通过；失败日志保留。\\
2 & 四维 Wilson flow & RK3 步长收敛；plaquette 单调诊断与流后有限性通过。\\
3 & flowed Clover、对偶场强与 Wilson 线 & traceless/legacy 投影明确；$\pm z$ 与 staple 几何合法；$z=0$ odd 为零。\\
4 & 六组件 OPE 与 schema 产物 & NPY/JSON 原子配对；轴、dtype、哈希和 \texttt{not\_applied} 清单完整。\\
5 & 两点函数与三点函数 & 源汇、宇称投影和真空扣除固定；$t_{\mathrm{sep}}$ 扫描有端点切割。\\
6 & ratio-v2 与 delete-one jackknife & 共同 nopol 分母、复数协方差、NaN 掩码和 selector 均通过。\\
7 & 基态矩阵元 $h(z,\bperp,\Pz)$ & plateau/多态拟合在预注册窗口内稳定；不按点挑窗口。\\
8 & 重整化方案 & 明确选择式~\eqref{eq:b10-chain-gf} 或式~\eqref{eq:b10-chain-hybrid}，方案因子不混用。\\
9 & Fourier、CS、软函数与 NLO 匹配 & 同一表示/几何/快度方案；$z_{\max}$ 与尾部模型误差可见。\\
10 & 联合外推与物理发布 & $a,\Pz,m_\pi,L$ 的窗口和相关性受控；所有未闭合项写入结果边界。\\
\bottomrule
\end{longtable}

\section{ratio-v2：共同分母、复数相位与掩码}

对每个通道 $\chi$，三点估计量应由配置/源协方差定义，而不是先对各组态取比值。
本批次新版采用的写法是
\begin{align}
  C_{3,\chi}
  &=\langle O_\chi C_{2,\chi}\rangle_{\mathrm{conf,src}}
  -\langle O_\chi\rangle_{\mathrm{conf,src}}
  \langle C_{2,\chi}\rangle_{\mathrm{conf,src}},
  \label{eq:b10-c3}\\
  C_{2,\mathrm{den}}
  &=\langle C_{2,\mathrm{nopol}}\rangle_{\mathrm{conf,src}},
  \qquad
  R_\chi=\frac{C_{3,\chi}}{C_{2,\mathrm{den}}}.
  \label{eq:b10-ratio}
\end{align}
非极化和螺旋度的 numerator 可以使用各自的通道，但分母必须共享 nopol。
尤其不能把 ensemble mean 的 pol35 作为近零分母。对螺旋度，完整复数必须先
经过协方差、真空扣除、ratio 和 jackknife，最后再做
\begin{equation}
  \operatorname{Re}[-\mathrm{i}R_H]=\operatorname{Im}R_H.
  \label{eq:b10-phase}
\end{equation}
若提前把 pol35 截成实部或虚部，式~\eqref{eq:b10-c3} 中的协方差项已经被
改变，最后再取虚部不能恢复。

delete-one jackknife 的每一样本应重新计算分子、共同分母和真空扣除：
\begin{equation}
  R_{(-k)}
  =\frac{C_{3,(-k)}}{C_{2,(-k),\mathrm{den}}},
  \qquad
  \sigma_{\mathrm{Re}}^2=\frac{N-1}{N}\sum_k
  \left(\operatorname{Re}R_{(-k)}
  -\overline{\operatorname{Re}R}\right)^2,
  \label{eq:b10-jk}
\end{equation}
虚部误差同理，实部与虚部不可合并成没有定义的单一标量。无效插入位置
$0\le\tau_{\mathrm{ins}}\le t_{\mathrm{sep}}$ 之外必须填 NaN，并保存
\texttt{valid\_insertion\_mask}。填零会把“没有观测”误当作“观测到零”，
从而系统性拉偏均值和误差。

\section{可复现验证与不能越过的边界}

新版工作流在来源中登记了三级证据：提交/标签说明、九项独立契约和全量回归。
其中可登记的断言包括：复杂 ratio 与参考估计器逐项一致、端点转换与参考一致、
随机局域规范变换下双局域 trace 不变、RK3 细步长收敛，以及全量回归通过。
这些数字是 U077 快照内的验证声明；B10 只做了文档和 PDF 元数据核对，没有
执行该快照中的安装、GPU、生产数据处理或测试命令。因此本草案把它们写成
“来源内登记证据”，而不是本次复跑结果。

同样的边界适用于物理结果。有限流时间的 flowed bare OPE 或 bare ratio 是
有意义的中间可观测量，但还不是物理胶子 PDF/TMD。当前材料明确把以下对象
列为未闭合或需真实生产接线：参考组态共同 manifest、单态 gluon--quark
mixing、与 schema-v2 ratio 相接的软函数/快度减除、完整 NLO 匹配、多格距
连续极限和场强树图归一 $\kappa_F$。只有在这些条件同时满足后，才可写出带
统计和系统误差的 $xg(x,\bperp;\mu,\zeta)$。

\section{历史谱系：从 RI/MOM 到梯度流 TMD}

本批次资料呈现了清晰但不等权的演进链。2022 年硕士学位论文以准 PDF 的 UV
线性发散为中心，摘要报告 RI/MOM 不能在所有格距上完全消除 Wilson 线发散；
2023 年硕士学位论文把 TMD-PDF、Wilson loop、HMC/Chroma 与 jackknife/
bootstrap 组织成格点工作流，但本批次只核对了题录和目录。两者提供历史问题
与工作流线索，不能替代 U077 中的显式公式。

旧版 \texttt{理论解析与工作流.tex} 采用共线胶子 PDF 主线，以 Eq.(20) 的
矩阵元、Eq.(25) 的非极化 OPE、$C_3/C_2$、Fourier 和 NLO matching 组织全文。
新版 \texttt{理论解析与工作流-v20260911.tex} 明确把它标为过时版本，并将主线
升级为梯度流重整化的胶子 TMD-PDF：新增 `SFTE` 端点转换、ratio-v2 复数统计、
CS 软结构、schema-v2 和验证边界。工作流版本因此应去重，而不是把两套公式
并列成两个物理结论。

从方法史看，可把演进概括为四步：先由 twist-2 定域算符计算低阶 Mellin 矩；
再由 LaMET/quasi-PDF 获得 $x$ 依赖；随后把直线 Wilson 线推广到 staple 并加入
软函数、CS 核和 TMD 因子化；最后用梯度流、系统误差窗口和多格距外推控制 UV
与连续极限。这个链条的每一步都增加新的标度和新的失效条件；后一步不会自动
修正前一步的错误。

\section{系统误差与失败模式}

\begin{tabularx}{\textwidth}{@{}p{0.22\textwidth}X@{}}
\toprule
来源 & 物理含义与最低报告要求 \\
\midrule
离散化 $a^2$ & Clover 改进后仍需至少三格距或明确说明只有单格距；流化半径必须
沿固定物理序列变化。 \\
流时间 SFTE & 同时检查式~\eqref{eq:b10-sfte-window}；小流截距与大流诊断分开，
不能以统计误差最小替代窗口控制。 \\
有限体积与激发态 & 报告 $M_\pi L$、$t_{\mathrm{sep}}$、端点切割和多态拟合；
disconnected 近似不自动消除所有污染。 \\
Fourier 截断 & 报告 $z_{\max}$、大 $z$ 尾部模型、Gibbs 振荡和逆问题正则化；
负值不能一律归因于统计。 \\
软函数/快度 & 报告 $S_I$、$\zeta$、Wilson 线方向与表示；缺少软减除时不能把
准 TMD 称为物理 TMD。 \\
匹配与混合 & 报告 $NLO$ 阶数、$Z_{ij}$、负半轴、$\mu$ 变异和单态混合；
正半轴截断或方案混用会造成不可忽略的伪偏移。 \\
方案与归一 & 固定 GF 或 hybrid 链；记录 $\kappa_F$、场强投影、Clover/对偶
约定和 metadata 中的未应用项。 \\
\bottomrule
\end{tabularx}
